\section{Executable validation, theorem status, and reproducibility}

\subsection{Purpose of the executable companion}

The file \code{code/zariski\_descent.py} is an exact-arithmetic reference implementation of the
central finite algorithms.  It uses only the Python standard library and never uses floating point.
It is not an iPRA proof checker and does not replace the universal arguments in the preceding
sections.  Its role is to check formulas, sign conventions, exponent bookkeeping, and examples
independently of the LaTeX derivation.

The implementation includes:
\begin{itemize}
\item a minimal exact commutative-ring interface, with instances for $\Z$ and $\Z/m\Z$;
\item verification and exponentiation of unit-ideal cover certificates;
\item the multinomial power-cover coefficient constructor;
\item principal-cover gluing for the module $M=A$;
\item sparse univariate polynomial arithmetic and ideal-certificate checking;
\item the $\mathbf P^1$ Laurent \v{C}ech differential and normal form;
\item projective-space cohomology rank formulas;
\item integer matrices for the monomial-wise differential and insertion homotopy.
\end{itemize}

\subsection{Recorded validation}

Running
\begin{lstlisting}[language={}] 
python3 -m unittest discover -s code -v
python3 code/zariski_descent.py
\end{lstlisting}
produced the log in \code{verification/python\_tests.txt} and the structured records in
\code{metadata/validation.json} and \code{verification/demo\_output.json}.  The eight unit tests
passed.  The recorded checks comprise:

\begin{longtable}{@{}p{0.32\textwidth}p{0.58\textwidth}@{}}
\toprule
Component & Recorded check \\
\midrule
\endfirsthead
\toprule
Component & Recorded check \\
\midrule
\endhead
Power covers & Eleven exact cases over $\Z$ and $\Z/12\Z$, including exponents larger than one. \\
Integer gluing & A two-chart example with unequal denominators recovers the global section $17$; the optimized common exponent is $2$, and the computed power-cover coefficients are $(7,-3)$. \\
Zero-divisor gluing & A cover in $\Z/12\Z$ recovers the global section $5$, demonstrating that the algorithm does not cancel nonunits or assume a domain. \\
Ideal certificate & An exact polynomial identity of the form $p-q=hg$ is accepted coefficientwise; a perturbed certificate is rejected. \\
$\mathbf P^1$ normal form & $1{,}200$ randomly generated Laurent cocycles satisfy exact reconstruction, and every obstruction term lies in $d<e<0$. \\
$\mathbf P^1$ global sections & Generated coefficient vectors give compatible chart pairs for several positive degrees and the expected empty basis for negative degrees. \\
Projective ranks & The code checks the closed rank formulas against direct bounded enumeration in the tested range. \\
$\mathbf P^n$ homotopy & For all proper nonempty negative-coordinate masks in dimensions $1\leq n\leq6$, $240$ matrix cases satisfy $dh+hd=I$ over $\Z$. \\
\bottomrule
\end{longtable}

The report uses the uniformly safe gluing exponent $q=N+K+1$.  The executable companion uses the
slightly sharper $q=\max(1,N+K)$ and then verifies the resulting cross identities.  This difference
is intentional: the paper favors a simple bound with slack, while the program demonstrates that the
slack is not mathematically essential.

\subsection{Worked principal-cover examples}

\paragraph{An integer example.}
Take the cover elements $f_0=2$, $f_1=3$ with certificate
\[
                         (-1)\cdot2+1\cdot3=1.
\]
Choose local fractions obtained from the global integer $17$ but with different denominator
exponents.  The compatibility checker clears the denominators, the gluing routine constructs a
common power cover, and the output is exactly $17$.  This test exercises denominator homogenization
rather than merely gluing two identical numerators.

\paragraph{A zero-divisor example.}
In $A=\Z/12\Z$, the elements $3$ and $4$ satisfy
\[
                         3\cdot3+1\cdot4=1\pmod{12}.
\]
Neither chart element is a unit, and $A$ has zero divisors.  Localizations and overlap equalities are
therefore genuinely proof-relevant.  The supplied local data glue to $5\in A$.  Every final
restriction equation is checked in the ring.  This example is a useful guard against accidentally
using cancellation in the human or Lean proof.

\subsection{Claim ledger}

The machine-readable file \code{metadata/claims.json} and the human-readable
\code{metadata/THEOREM\_LEDGER.md} distinguish the four relevant levels of completion.  The main
status is summarized here.

\begin{longtable}{@{}p{0.12\textwidth}p{0.48\textwidth}p{0.31\textwidth}@{}}
\toprule
ID & Result & Status in this archive \\
\midrule
\endfirsthead
\toprule
ID & Result & Status in this archive \\
\midrule
\endhead
F1 & Certified-setoid and congruence discipline & Full human proof and Lean interface; object derivation pending. \\
A1 & Quotient equality by ideal combinations & Full human proof and exact checker; object derivation pending. \\
A2 & Localization equivalence and ring congruence & Explicit witness formulas proved; Lean target supplied. \\
Z1 & Power-cover constructor & Full proof, exact implementation, finite tests; \code{Prf PRAAx} pending. \\
Z2 & Joint monicity on a principal cover & Full proof; object derivation pending. \\
Z3 & Principal-cover module gluing & Full proof including zero divisors; exact implementation for $M=A$; object derivation pending. \\
G1--G2 & Finite affine atlases and quasi-coherent descent data & Precise certificate semantics and constructions; no complete category-of-schemes library claimed. \\
P1--P3 & $\mathbf P^1$, $\cO(d)$, $H^0$, and $H^1$ & Full finite proofs and exact normal-form tests; object derivations pending. \\
PN & Mixed-sign contraction and the standard $\mathbf P^n$ table & Full combinatorial proof and finite matrix checks; object derivation pending. \\
C1 & Agreement with conventional sheaf cohomology & Uses the standard affine-cover comparison theorem; no internal derived-functor formalization claimed. \\
L1 & Lean completion route & Source-level blueprints and dependency order; not compiler-checked. \\
\bottomrule
\end{longtable}

\subsection{Negative checks and nonclaims}

The validation deliberately does not infer an unbounded theorem from a successful finite test.
Specifically:
\begin{itemize}
\item testing finitely many rings does not prove the power-cover theorem for every presented ring;
\item a bounded point search does not prove absence of points above its cutoff;
\item a successful certificate checker does not guarantee that a certificate producer exists for
      every true semantic statement;
\item native Lean propositions, once added, would not by themselves be \code{Prf PRAAx} values;
\item the standard \v{C}ech computation is not by itself a formalization of all derived-functor
      sheaf cohomology inside iPRA.
\end{itemize}

No claim is made here of general effective Siegel or Faltings, uniform Mordell--Weil generators,
exact singular Puiseux selection from arbitrary Cauchy coefficients, decidable equality of arbitrary
represented reals, or a primitive-recursive universal evaluator.

\subsection{Rebuilding the package}

From the extracted package root, the intended commands are:
\begin{lstlisting}[language={}] 
python3 -m unittest discover -s code -v
python3 code/zariski_descent.py > metadata/validation.json
cd paper
latexmk -pdf -interaction=nonstopmode -halt-on-error \
  ipra_algebraic_geometry_framework.tex
\end{lstlisting}
The included \code{build.sh} performs the tests, rebuilds the PDF, and writes an integrity manifest.
A normal TeX Live installation with the packages named in the master source is sufficient.

The upstream repository audit is pinned to
\code{bce5113f44d613f5625c274128ee761f35486fc5}.  At the final check, \code{main} was
\code{a2f72b0cdf57f8a3fefd34b272c3d471385e819c}; the Git object identifiers of \code{ConciseIPRA}, the iPRA
\code{README.md}, and \code{docs} were identical at the two commits.  The archive neither
modifies nor rebuilds the repository.  Its source audit is in
\code{metadata/SOURCE\_AUDIT.md}; its explicit limitations are in
\code{metadata/LIMITATIONS.md}.

\subsection{Archive inventory}

\begin{longtable}{@{}p{0.35\textwidth}p{0.55\textwidth}@{}}
\toprule
Path & Purpose \\
\midrule
\endfirsthead
\toprule
Path & Purpose \\
\midrule
\endhead
\code{paper/ipra\_algebraic\_geometry\_framework.tex} & Master editable LaTeX source. \\
\code{paper/sections/} & The complete section sources, including proofs and formalization plan. \\
\code{paper/ipra\_algebraic\_geometry\_framework.pdf} & Compiled report. \\
\code{code/zariski\_descent.py} & Exact finite reference algorithms. \\
\code{code/test\_zariski\_descent.py} & Unit and randomized regression tests. \\
\code{lean/*.lean} & Host structures, target propositions, and object-theory endpoint blueprint. \\
\code{lean/FORMALIZATION\_PLAN.md} & Dependency-ordered path to complete Lean and object-iPRA proofs. \\
\code{metadata/SOURCE\_AUDIT.md} & Pinned repository and external-source audit. \\
\code{metadata/THEOREM\_LEDGER.md} & Result-by-result proof status. \\
\code{metadata/claims.json} & Machine-readable claim levels and nonclaims. \\
\code{metadata/validation.json} & Exact executable validation record. \\
\code{verification/} & Test logs, demo output, PDF build report, and visual QA notes. \\
\code{MANIFEST.sha256} & Integrity hashes of delivered files other than the manifest itself. \\
\bottomrule
\end{longtable}
